\subsection{Consolidação do planejamento complementar e contrato de execução}
\label{sec:consolidacao}
O inventário complementar inclui Planejamento\_Etiquetas\_LCQUI.md, REGRAS\_MULTI\_ROLE.md, COMPILACAO\_NIX\_LCQUI.md, README.md e os quatro documentos de acompanhamento. Os README de scaffold do frontend não adicionam requisitos de negócio; instruções AGENTS/CLAUDE orientam desenvolvimento, não a especificação do produto. DUVIDAS\_PENDENTES\_LCQUI.md é o registro atual de decisões; Q01--Q14 e DP-A01--DP-D02 já foram resolvidas, conforme MODIFICACOES\_CONSOLIDADAS\_LCQUI.md.

O planejamento de etiquetas foi incorporado às seções 4 (auditoria), 5.9 (contador e campos), 7 (regras), 8 (UI-08), 9 (ALM-07), 10 (geração) e 11 (acesso). O documento de Multi-Role foi incorporado integralmente pelos IDs RN-ROLE-01 a RN-ROLE-15 na seção 7. Os Markdown de origem permanecem como histórico; esta versão explicita as decisões conflitantes em vez de criar fontes concorrentes.

\paragraph{Etiquetas e geometria.}
Papel A4 210 por 297 mm; três colunas e dez linhas; etiqueta 65 por 26,5 mm; margens laterais 7,5 mm e superior/inferior 16 mm; espaçamentos horizontal/vertical zero. Linhas cinza contínuas orientam corte. Offset é aplicado apenas na primeira página; ao completar a grade, reiniciar na célula 1/1 da seguinte. Etiqueta virgem contém código, espaço manuscrito para data/reagente e Code 128. Segunda via contém nome e data cadastrados e uma etiqueta centralizada por ficha. Gerar barras em memória com bwip-js e PDF com pdfkit; preservar proporção e área livre ao redor do código, testar leitura física e códigos longos. Casos mínimos: uma etiqueta, coluna 3, offset 4/2, última célula 10/3, 30/36 itens, nome longo, frasco sem histórico e rejeição de 11 IDs na segunda via. Os testes relatados no planejamento são evidência histórica de protótipo, não homologação da aplicação atual.

\paragraph{Contador e auditoria.}
Adota-se o caminho já usado pelo sistema, Contador\_Codigo\_Frasco/singleton e ultimo\_codigo\_gerado. Este é o único caminho normativo do sequenciador, sem contador alternativo. Cadastro incrementa contador e cria frasco atomicamente; impressão não reserva código. Lotes virgens usam Impressao\_Etiqueta\_Frasco; segunda via usa Registro\_de\_Auditoria. A sugestão contraditória de armazenar segunda via também em Impressao não é adotada. Custos reais incluem todas as escritas/leitura de domínio, índices, histórico e tentativas; não reutilizar estimativa de ``1 leitura e 2 escritas'' como custo total.

\paragraph{Relatórios, transporte e integridade.}
O backend atual retorna PDF em base64; a especificação adota esse transporte para relatórios limitados, com conversão em Blob no cliente e liberação da URL temporária. Roteiros e comprovantes permanecem no Storage. Relatórios maiores exigem decisão de infraestrutura antes de ampliar limites. Um hash simples de parâmetros/data não é assinatura digital nem prova de autoria; tratá-lo como identificador de rastreabilidade. Relatórios históricos usam eventos/snapshots da época, não somente o cadastro atual.

\paragraph{Garantias de execução.}
Exemplos de código desta seção são ilustrativos e não substituem contratos UI/dicionário. Todas as leituras transacionais precedem escritas; callbacks podem repetir e não devem enviar e-mail ou gerar PDF como efeito externo. A transação grava evento/operação e um processador idempotente realiza efeitos externos após commit (contrato global de idempotência na Seção~7). Alterações no Auth e upload no Storage exigem reconciliação, pois não participam da transação Firestore. Consultas in e getAll não têm o mesmo contrato: tamanho de lote deve ser escolhido por API e medido; 200 não é limite universal do Firestore.

\paragraph{Ambiente e documentação operacional.}
README contém instruções de aplicação e emuladores; COMPILACAO\_NIX\_LCQUI.md é guia de build documental. O flake.nix atual fornece Node/Python/Git/Docker, mas não TeX Live; não presumir que nix develop disponibiliza pdflatex. Usar TeX Live instalado ou ambiente Nix dedicado e compilar main.tex até estabilizar referências. Os quatro documentos de acompanhamento registram evidências do código, não novos requisitos independentes.

\paragraph{Referências técnicas verificadas.}
\href{https://firebase.google.com/docs/firestore/manage-data/transactions}{Firebase: transações e lotes};
\href{https://firebase.google.com/docs/auth/admin/custom-claims}{Firebase: propagação de Custom Claims};
\href{https://firebase.google.com/docs/firestore/manage-data/enable-offline}{Firebase: persistência offline}.
Consulta em 11/09/2026. Estes fundamentos apoiam a execução transacional, renovação de tokens e opção por cache em dispositivo confiável.

\paragraph{Contratos dos auxiliares Q04/Q14.}
validarAceiteTcrTx lê o evento de aceite produzido em sessão autenticada do retirante, confere versão institucional vigente, UID, frasco, finalidade, operação ainda não consumida e justificativa de 20--2000 caracteres, verificando o mínimo após trim por validarJustificativaHumana (Seção 10.5). Não aceita timestamp ou UID arbitrário no payload do gestor. vincularAceiteEAuditoriaTx consome o aceite e registra o empréstimo correspondente na mesma transação. validarAutoAtendimentoTx conta gestores ativos, bloqueia quando existe outro e exige justificativa humana de pelo menos 20 caracteres após trim, usando validarJustificativaHumana quando a condição de autoatendimento se aplica; registrarNotificacaoChefiaTx grava notificação/outbox idempotente. resolverSnapshotQuimicoTx resolve unidade pelo resumo e densidade pela especificação, sem gravações antes de concluir leituras. Auxiliares ilustram contratos, não atestam implementação em functions/. O TCR não substitui segurança química institucional.

\subsubsection{Contratos complementares pós-auditoria}

\begin{lstlisting}[language=TypeScript, title={M-12: parseDataCivil -- datas civis normalizadas em America/Sao\_Paulo}]
// M-12: todas as datas civis YYYY-MM-DD devem ser interpretadas no fuso
// America/Sao_Paulo, não em UTC. `new Date("2026-09-15")` retorna meia-noite UTC
// (=21:00 do dia anterior em Sao_Paulo), o que causaria comparações erradas.
// Prazo de devolução encerra-se às 23:59:59.999 do dia previsto em Sao_Paulo.
//
// Nota: horário de verão histórico (Brasil aboliu em 2019) não é mais risco futuro,
// mas o helper usa Intl para robustez caso o fuso mude.
import { DateTime } from "luxon";

// Interpreta como data local em America/Sao_Paulo às 00:00:00
function parseDataCivil(dataStr: string): Date {
  return DateTime.fromISO(dataStr, { zone: "America/Sao_Paulo" }).startOf("day").toJSDate();
}

function prazoLimiteDevolucao(dataStr: string): Date {
  return DateTime.fromISO(dataStr, { zone: "America/Sao_Paulo" }).endOf("day").toJSDate();
}
\end{lstlisting}

\begin{lstlisting}[language=TypeScript, title={M-11: geração de codigo\_turma -- base32 Crockford 6 chars com retry transacional}]
// M-11: codigo_turma UNIQUE, <= 20 chars, gerado pelo servidor.
// Usa Crockford Base32 (alfabeto sem I, L, O, U) para legibilidade.
// Retry transacional via docId determinístico em coleção Chaves_Unicas.
const CROCKFORD = "0123456789ABCDEFGHJKMNPQRSTVWXYZ";

function gerarCodigoTurmaCandidate(): string {
  const bytes = require("crypto").randomBytes(4); // 32 bits => 6 chars Crockford
  let n = bytes.readUInt32BE(0);
  let codigo = "";
  for (let i = 0; i < 6; i++) {
    codigo = CROCKFORD[n & 0x1f] + codigo;
    n >>= 5;
  }
  return codigo;
}

async function gerarCodigoTurmaTx(
  tx: FirebaseFirestore.Transaction
): Promise<string> {
  for (let tentativa = 0; tentativa < 5; tentativa++) {
    const candidato = gerarCodigoTurmaCandidate();
    const lockRef = admin.firestore()
      .collection("Chaves_Unicas").doc(`Turma__${candidato}`);
    const lockSnap = await tx.get(lockRef);
    if (!lockSnap.exists) {
      tx.create(lockRef, { reservado_em: admin.firestore.FieldValue.serverTimestamp() });
      return candidato;
    }
  }
  throw new HttpsError("internal", "Falha ao gerar codigo_turma único após 5 tentativas.");
}
\end{lstlisting}


\begin{lstlisting}[language=TypeScript, title={N-12: aceitarConviteAluno com autenticação prévia e criação condicional Firestore}]
// N-12: quem aceita o convite já está autenticado em Firebase Auth
// (request.auth.uid presente) e com e-mail verificado; a conta Auth do
// solicitante já existe (cadastro/login prévio). Criação condicional refere-se
// apenas a Usuarios/Aluno no Firestore, reutilizando o UID existente pelo
// e-mail para evitar identidade duplicada. O papel Aluno e o vínculo são
// verificados/criados na mesma transação Firestore; nenhuma chamada externa
// (Auth/e-mail) ocorre dentro da transação. O retry idempotente do mesmo ator
// devolve o receipt; criação de outra conta Auth não faz parte deste fluxo.
export const aceitarConviteAluno = onCall(async (request) => {
  const { tokenConvite, matriculaInformada } = request.data as { tokenConvite: string; matriculaInformada?: string };
  if (!request.auth?.uid) throw new HttpsError("unauthenticated", "Usuario nao autenticado.");

  // 1. Calcular hash SHA-256 para busca do convite pelo token_hash (o HMAC e usado apenas como ID deterministico do documento)
  const { createHash } = require("crypto");
  const hashToken = createHash("sha256").update(tokenConvite).digest("hex");

  // 2. Verificar e-mail do usuario autenticado no provedor
  const authUser = await admin.auth().getUser(request.auth.uid);
  const emailNorm = authUser.email?.toLowerCase();

  // 3. Transacao: revalidar convite e criar Aluno e vinculos se nao existirem
  return admin.firestore().runTransaction(async (tx) => {
    // PASSO A: Todas as leituras
    const conviteQuery = await tx.get(admin.firestore().collection("Convite_Aluno")
      .where("token_hash", "==", hashToken).limit(1));
    if (conviteQuery.empty) throw new HttpsError("not-found", "Convite invalido.");

    const conviteDoc = conviteQuery.docs[0];
    const convite = conviteDoc.data();

    if (convite.status !== "pendente")
      throw new HttpsError("failed-precondition", "Convite ja aceito ou expirado.");
    if (new Date() > convite.expira_em.toDate())
      throw new HttpsError("failed-precondition", "Convite expirado.");
    if (emailNorm !== convite.email)
      throw new HttpsError("permission-denied", "E-mail da conta nao corresponde ao convite.");

    const uid = request.auth!.uid;
    const alunoRef = admin.firestore().collection("Aluno").doc(uid);
    const usuarioRef = admin.firestore().collection("Usuarios").doc(uid);

    const alunoSnap = await tx.get(alunoRef);
    const usuarioSnap = await tx.get(usuarioRef);

    let turmaSnap = null;
    let turma = null;
    let jaMatriculado = false;
    let alunosSnap = null;

    if (convite.id_turma) {
      const turmaRef = admin.firestore().collection("Turma").doc(convite.id_turma);
      turmaSnap = await tx.get(turmaRef);
      if (!turmaSnap.exists) throw new HttpsError("not-found", "Turma nao encontrada.");
      turma = turmaSnap.data()!;
      if (turma.status !== "Ativo") throw new HttpsError("failed-precondition", "Turma inativa.");

      const alunoNaTurmaSnap = await tx.get(turmaRef.collection("Alunos").doc(uid));
      jaMatriculado = alunoNaTurmaSnap.exists;

      // Validacao de capacidade transacional (PDF-008: evita leitura O(N))
      if (!jaMatriculado && !convite.exceder_capacidade) {
        const qtdAlunosAtual = turma.qtd_alunos || 0;
        if (qtdAlunosAtual >= turma.capacidade) {
          throw new HttpsError("failed-precondition", "Turma lotada.");
        }
      }
    }

    // PDF-003: Validar matricula e unicidade se for criar um novo Aluno
    let lockRefMatricula: FirebaseFirestore.DocumentReference | null = null;
    let matriculaFinal: string | null = null;
    if (!alunoSnap.exists) {
      const matriculaBruta = convite.numero_matricula ?? matriculaInformada;
      if (typeof matriculaBruta !== "string") {
        throw new HttpsError("invalid-argument", "Matrícula é obrigatória e deve ser texto.");
      }

      matriculaFinal = matriculaBruta.trim(); // Canonicalização documental

      if (!matriculaFinal) {
        throw new HttpsError("invalid-argument", "Matrícula não pode ser vazia ou conter apenas espaços.");
      }
      if (matriculaFinal.length > 20) {
        throw new HttpsError("invalid-argument", "Matrícula inválida: excede o limite de 20 caracteres.");
      }

      lockRefMatricula = admin.firestore().collection("Chaves_Unicas").doc(`Aluno__${matriculaFinal}`);
      const lockSnap = await tx.get(lockRefMatricula);
      if (lockSnap.exists) throw new HttpsError("already-exists", "A matrícula informada já está em uso por outro aluno.");
    }

    // PASSO B: Escritas (idempotentes para repeticao concorrente)
    if (!usuarioSnap.exists) {
      tx.set(usuarioRef, {
        nome: authUser.displayName ?? emailNorm,
        email: emailNorm,
        ativo: true,
        profile_url_photo: authUser.photoURL ?? null,
      });
    }
    if (!alunoSnap.exists && matriculaFinal && lockRefMatricula) {
      tx.set(alunoRef, {
        id_usuario: uid,
        numero_matricula: matriculaFinal,
      });
      tx.set(lockRefMatricula, {
        tipo: "Aluno",
        id_recurso: uid,
        criado_em: admin.firestore.FieldValue.serverTimestamp()
      });
    }

    // Expirar convite e reconciliar com Secao 5 (aceitado_por, aceitado_em)
    const agora = admin.firestore.FieldValue.serverTimestamp();
    tx.update(conviteDoc.ref, {
      status: "aceitado",
      aceitado_por: uid,
      aceitado_em: agora
    });

    // Vincular a turma (se aplicavel e se nao estiver matriculado ainda)
    if (convite.id_turma && turmaSnap && !jaMatriculado) {
      const turmaRef = turmaSnap.ref;
      const vincRef = turmaRef.collection("Alunos").doc(uid);
      const vincRef2 = usuarioRef.collection("Turmas").doc(convite.id_turma);

      // 1o Espelho: colecao Alunos na Turma
      tx.set(vincRef, { id_aluno: uid, ingressou_em: agora });
      // 2o Espelho: colecao Turmas no Usuario
      tx.set(vincRef2, {
        id_turma: convite.id_turma,
        id_professor: turma.id_professor,
        id_materia: turma.id_materia,
        nome_turma: turma.nome_turma,
        nome_materia: turma.nome_materia,
        ano: turma.ano,
        semestre: turma.semestre,
        status: turma.status,
        ingressou_em: agora
      });
      // Registro no historico: preserva regras de reingresso sem duplicacao
      tx.set(turmaRef.collection("HistoricoAlunos").doc(), {
        id_aluno: uid,
        tipo: "inclusao_aluno",
        modo_ingresso: "CONVITE",
        justificativa: convite.justificativa_excecao || null,
        timestamp: agora
      });
      tx.update(turmaRef, { qtd_alunos: admin.firestore.FieldValue.increment(1) });
    }

    return { vinculado: !!convite.id_turma, ja_matriculado: jaMatriculado };
  });
  // Apos commit: atualizarCustomClaims(request.auth!.uid) -- recuperavel e idempotente
});
\end{lstlisting}

\paragraph{Contrato M11 (turmas, matrícula e convites).}
O pseudocódigo N-12 é ilustrativo. O contrato normativo é a
Seção~\ref{sec:regras-turmas-m11}: o vínculo canônico é
\texttt{Turma/\{id\}/Alunos/\{uid\}} e o espelho
\texttt{Usuarios/\{uid\}/Turmas/\{id\}} é projeção; \texttt{qtd\_alunos} é lido
e escrito na transação, nunca recontado fora dela; o docId do convite é
determinístico por HMAC do contexto e do e-mail normalizado, com unicidade em
\texttt{Chaves\_Unicas}; o ID do convite é imutável e distinto da chave de
pendência, de modo que um novo convite após terminalidade não sobrescreve o
aceito/expirado nem apaga \texttt{aceitado\_por}/\texttt{aceitado\_em}; a
aceitação exige e-mail verificado correspondente e é idempotente por (ator,
token, operação) M7, distinguindo retry de tentativa distinta; convite global
não cria matrícula; a criação condicional de \texttt{Usuarios}/\texttt{Aluno} é
Firestore e transacional, enquanto o envio de e-mail e a sincronização de
Custom Claims são etapas externas, pós-commit, recuperáveis e idempotentes;
criar o registro não é enviar o e-mail. M-11 reserva \texttt{codigo\_turma} em
\texttt{Chaves\_Unicas} e arquivar não libera o código. O fan-out de
\texttt{alterarStatusTurma} para os espelhos é idempotente e não é autoridade
de escrita (Q08).

\begin{lstlisting}[language=TypeScript, title={M-01: registrarBaixaBemPatrimonial (RF12)}]
// M-01: Regra RF12 -- registro de baixa patrimonial após processo institucional.
// O gestor faz upload do PDF comprobatório antes de chamar esta função.
// O backend valida a existência do documento no Storage antes de confirmar a baixa.
export const registrarBaixaBemPatrimonial = onCall(async (request) => {
  const { idBem, documentoUrl } = request.data as { idBem: string; documentoUrl: string };
  await validarPermissao(request, ["Chefe_Geral", "Gestor_Bens_Patrimoniais"], true);
  await validarDocumentoStorage(documentoUrl, request.auth!.uid, "PDF"); // existência, ACL, tamanho e magic bytes %PDF-; fixa geração validada
  return admin.firestore().runTransaction(async (tx) => {
    const bemRef = admin.firestore().collection("Bem_Patrimonial").doc(idBem);
    const snap = await tx.get(bemRef);
    if (!snap.exists) throw new HttpsError("not-found", "Bem não encontrado.");
    const bem = snap.data()!;
    if (bem.status !== "Inservivel")
      throw new HttpsError("failed-precondition", "Bem deve estar Inservível antes da baixa.");
    tx.update(bemRef, {
      status: "Ja_dado_baixa",
      documento_dado_baixa_pdf_url: documentoUrl,
      versao: bem.versao + 1,
    });
    registrarHistoricoBemTx(tx, bemRef, bem, { status: "Ja_dado_baixa" }, request.auth!.uid);
    return { status: "Ja_dado_baixa" };
  });
});
\end{lstlisting}

\begin{regra}[M-07: Bounds e validação na geração de etiquetas]
O grid de etiquetas é sempre 3 colunas × 10 linhas (30 posições por página A4). A seleção de offset inicial aplica-se apenas na primeira página, validando: coluna em $[1, 3]$ e linha em $[1, 10]$. Lotes de até 50 etiquetas virgens são suportados por operação; as páginas subsequentes preenchem as posições continuamente desde a coordenada (1,1). Para segunda via, o limite é de 10 frascos por sessão, validado no backend antes de iniciar a geração. IDs fora do intervalo de etiquetas virgens conhecidas são rejeitados com mensagem explícita.
\end{regra}

